% \begin{frame}{Validation Settings \& Experimental Conditions}
% \framesubtitle{Testing weaker vs. stronger construct representations across two domains}

% \begin{columns}[t]
%     % LEFT COLUMN: FeedbackQA
%     \begin{column}{0.46\textwidth}
%         {\large \textbf{FeedbackQA}}\\[0.1cm]
%         \hrule
%         \vspace{0.2cm}
%         \textit{Public Q\&A dataset}\\[0.1cm]
%         Validity evidence: \textbf{Convergent}
%         (Alignment with human ratings)
        
%         \vspace{0.4cm}
%         \textbf{Specifications compared:}
%         \begin{itemize}
%             \item Naive holistic (1 score)
%             \item Informed holistic (1 score)
%             \item Formative criteria (4 scores)
%         \end{itemize}
%     \end{column}

%     % RIGHT COLUMN: Barter Deals
%     \begin{column}{0.46\textwidth}
%         {\large \textbf{Barter Deals}}\\[0.1cm]
%         \hrule
%         \vspace{0.2cm}
%         \textit{Proprietary marketplace data}\\[0.1cm]
%         Validity evidence: \textbf{Predictive} \\
%         (Alignment with application counts)
        
%         \vspace{0.4cm}
%         \textbf{Specifications compared:}
%         \begin{itemize}
%             \item Naive holistic (1 score)
%             \item Macro-formative (3 scores)
%             \item Fine-grained formative (10 scores)
%         \end{itemize}
%     \end{column}
% \end{columns}

% \vfill

% \begin{block}{Validation benchmarks}
% FeedbackQA compares LLM measurements against human ratings. Barter Deals tests whether LLM measurements add predictive value beyond a structural market baseline built from platform covariates.
% \end{block}

% \end{frame}


\begin{frame}{Validation Settings \& Experimental Conditions}
\framesubtitle{Testing LLM measurements against external benchmarks}

\begin{columns}[t]
    \begin{column}{0.46\textwidth}
        {\large \textbf{FeedbackQA}}\\[0.1cm]
        \hrule
        \vspace{0.2cm}
        \textit{Public Q\&A dataset}\\[0.1cm]
        Benchmark: \textbf{human ratings}

        \vspace{0.70cm}
        \textbf{Specifications compared:}
        \begin{itemize}
            \item Naive holistic (1 score)
            \item Informed holistic (1 score)
            \item Formative criteria (4 scores)
        \end{itemize}
    \end{column}

    \begin{column}{0.46\textwidth}
        {\large \textbf{Barter Deals}}\\[0.1cm]
        \hrule
        \vspace{0.2cm}
        \textit{Proprietary marketplace data}\\[0.1cm]
        Benchmark: \textbf{structural market baseline}\\
        Outcome: application counts

        \vspace{0.25cm}
        \textbf{Specifications compared:}
        \begin{itemize}
            \item Naive holistic (1 score)
            \item Macro-formative (3 scores)
            \item Fine-grained formative (10 scores)
        \end{itemize}
    \end{column}
\end{columns}

\vfill

\begin{block}{Key comparison}
Do explicit formative LLM measurements provide stronger validity evidence than weaker construct specifications or non-LLM baselines?
\end{block}

\end{frame}